\documentclass[11pt,a4paper]{article}
\begin{filecontents*}[overwrite]{tacl2021v1.sty}
\NeedsTeXFormat{LaTeX2e}
\ProvidesPackage{tacl2021v1}[2021/12/15 TACL formatting]
%
% tacl2021v1.sty: uses package options instead of a flag in the tex itself,
% and allows for a double-spaced version for use by copy-editors.
%
% Possible options are {hyperref,nohyperref}, {submission,acceptedWithA}, and {copyedit}
% Examples:
%
% \usepackage{tacl2021v1}   % equivalent to \usepackage[hyperref,submission]{tacl2021v1}
% \usepackage[acceptedWithA]{tacl2021v1}
% \usepackage[acceptedWithA,copyedit]{tacl2021v1}
%
% tacl2021v1.sty - (a) Updated the author/affiliation examples below to be
% consistent with the MIT Press-mandated options given in the new
% Appendix A of the template document. (b) updated the version number
% of "2021v1"
% -- Cindy Robinson <pvaql.genafnpy@tznvy.pbz> (<- ROT13)
%
% tacl2018.sty This is a modified version of ACL 2018 that replaces conference-specific
% information with TACL-specific information; removes some outdated comments;
% and corrects some minor spelling issues.
% --Lillian Lee (llee@cs.cornell.edu)
%
% ACL 2018: modified line numbering to be in gainsboro gray color so one
% can write comments on the margins.  -- Shay Cohen / Kevin Gimpel / Wei Lu

% 2017: modified to support DOI links in bibliography.  Now uses
% natbib package rather than defining citation commands in this file.
% Use with acl_natbib.bst bib style.  -- Dan Gildea

% This [was] the LaTeX style for ACL 2016. It contains Margaret Mitchell's
% line number adaptations (ported by Hai Zhao and Yannick Versley).
% It is nearly identical to the style files for ACL 2015,
% ACL 2014, EACL 2006, ACL2005, ACL 2002, ACL 2001, ACL 2000,
% EACL 95 and EACL 99.
% This [was] the LaTeX style file for ACL 2000.  It is nearly identical to the
% style files for EACL 95 and EACL 99.  Minor changes include editing the
% instructions to reflect use of \documentclass rather than \documentstyle
% and removing the white space before the title on the first page
% -- John Chen, June 29, 2000
% This [was] the LaTeX style file for EACL-95.  It is identical to the
% style file for ANLP '94 except that the margins are adjusted for A4
% paper.  -- abney 13 Dec 94
% The ANLP '94 style file is a slightly modified
% version of the style used for AAAI and IJCAI, using some changes
% prepared by Fernando Pereira and others and some minor changes
% by Paul Jacobs.
% Prior version prepared by Peter F. Patel-Schneider,

% Papers are to be prepared using the ``acl_natbib'' bibliography style,
% as follows:
%       \documentclass[11pt]{article}
%       \usepackage{tacl2021v1}
%       \title{Title}
%       \author{Author 1 \and Author 2 \\ Address line \\ Address line \And
%               Author 3 \\ Address line \\ Address line}
%       \begin{document}
%       ...
%       \bibliography{bibliography-file}
%       \bibliographystyle{acl_natbib}
%       \end{document}

% MIT Press mandates how the author block may be formatted.
% Please see the author block and appendix of the template file for
% instructions and examples.

% If the title and author information does not fit in the area allocated,
% place \setlength\titlebox{<new height>} right after
% \usepackage{tacl2021v1}
% where <new height> can be something larger than 5cm

%%%% hyperref should be the *last* package used %%%%

% Allow for double-spacing in copyedit version
\RequirePackage{setspace}

% include hyperref, unless user specifies nohyperref option like this:
% \usepackage[nohyperref]{tacl2018}
\newif\iftacl@hyperref
\DeclareOption{hyperref}{\tacl@hyperreftrue}
\DeclareOption{nohyperref}{\tacl@hyperreffalse}
\ExecuteOptions{hyperref} % default is to use hyperref

% non-anonymous and no ruler, vs anonymous with line-number ruler
\newif\iftaclpubformat % no @ so can use in instructions document.
\DeclareOption{acceptedWithA}{\taclpubformattrue}
\DeclareOption{acceptedwithA}{\taclpubformattrue} % allow for typos
\DeclareOption{acceptedwitha}{\taclpubformattrue} % allow for typos

\ExecuteOptions{submission} % default is submission version


% double spacing for copy-editor
\newif\iftaclcopyedit
\taclcopyeditfalse
\DeclareOption{copyedit}{\taclpubformattrue\taclcopyedittrue}

\ProcessOptions\relax


\iftacl@hyperref
  \RequirePackage{hyperref}
  \RequirePackage{xcolor}   % make links dark blue
  \definecolor{darkblue}{rgb}{0, 0, 0.5}
  \hypersetup{colorlinks=true,citecolor=darkblue, linkcolor=darkblue, urlcolor=darkblue}
\else
  % This definition is used if the hyperref package is not loaded.
  % It provides a backup, no-op definiton of \href.
  % This is necessary because \href command is used in the acl_natbib.bst file.
  \def\href#1#2{{#2}}
  % We still need to load xcolor in this case because the lighter line numbers require it. (SC/KG/WL)
  \RequirePackage{xcolor}
\fi





% Spacing idea communicated to L. Lee by Yoav Artzi
\newcommand{\Thanks}[1]{\thanks{\hspace{5pt} #1}}

% A4 modified by Eneko; again modified by Alexander for 5cm titlebox
\setlength{\paperwidth}{21cm}   % A4
\setlength{\paperheight}{29.7cm}% A4
\setlength\topmargin{-0.5cm}
\setlength\oddsidemargin{0cm}
\setlength\textheight{24.7cm}
\setlength\textwidth{16.0cm}
\setlength\columnsep{0.6cm}
\newlength\titlebox
\setlength\titlebox{5cm}
\setlength\headheight{5pt}
\setlength\headsep{0pt}
\thispagestyle{empty}
\pagestyle{empty}

\flushbottom \twocolumn \sloppy



% We're never going to need a table of contents, so just flush it to
% save space --- suggested by drstrip@sandia-2
\def\addcontentsline#1#2#3{}

%% ----- Set up hooks to repeat content on every page of the output doc,
%% necessary for the line numbers in the submitted version.  --MM
%%
%% Copied from CVPR 2015's cvpr_eso.sty, which appears to be largely copied from everyshi.sty.
%%
%% Original cvpr_eso.sty available at: http://www.pamitc.org/cvpr15/author_guidelines.php
%% Original evershi.sty available at: https://www.ctan.org/pkg/everyshi
%%
%% Copyright (C) 2001 Martin Schr\"oder:
%%
%%                         Martin Schr"oder
%%                         Cr"usemannallee 3
%%                         D-28213 Bremen
%%                         Martin.Schroeder@ACM.org
%%
%% This program may be redistributed and/or modified under the terms
%% of the LaTeX Project Public License, either version 1.0 of this
%% license, or (at your option) any later version.
%% The latest version of this license is in
%%    CTAN:macros/latex/base/lppl.txt.
%%
%% Happy users are requested to send [Martin] a postcard. :-)
%%
\newcommand{\@EveryShipoutTACL@Hook}{}
\newcommand{\@EveryShipoutTACL@AtNextHook}{}
\newcommand*{\EveryShipoutTACL}[1]
   {\g@addto@macro\@EveryShipoutTACL@Hook{#1}}
\newcommand*{\AtNextShipoutACL@}[1]
   {\g@addto@macro\@EveryShipoutTACL@AtNextHook{#1}}
\newcommand{\@EveryShipoutTACL@Shipout}{%
   \afterassignment\@EveryShipoutTACL@Test
   \global\setbox\@cclv= %
   }
\newcommand{\@EveryShipoutTACL@Test}{%
   \ifvoid\@cclv\relax
      \aftergroup\@EveryShipoutTACL@Output
   \else
      \@EveryShipoutTACL@Output
   \fi%
   }
\newcommand{\@EveryShipoutTACL@Output}{%
   \@EveryShipoutTACL@Hook%
   \@EveryShipoutTACL@AtNextHook%
      \gdef\@EveryShipoutTACL@AtNextHook{}%
   \@EveryShipoutTACL@Org@Shipout\box\@cclv%
   }
\newcommand{\@EveryShipoutTACL@Org@Shipout}{}
\newcommand*{\@EveryShipoutTACL@Init}{%
   \message{ABD: EveryShipout initializing macros}%
   \let\@EveryShipoutTACL@Org@Shipout\shipout
   \let\shipout\@EveryShipoutTACL@Shipout
   }
\AtBeginDocument{\@EveryShipoutTACL@Init}

%% ----- Set up for placing additional items into the submitted version --MM
%%
%% Based on eso-pic.sty
%%
%% Original available at: https://www.ctan.org/tex-archive/macros/latex/contrib/eso-pic
%% Copyright (C) 1998-2002 by Rolf Niepraschk <niepraschk@ptb.de>
%%
%% Which may be distributed and/or modified under the conditions of
%% the LaTeX Project Public License, either version 1.2 of this license
%% or (at your option) any later version.  The latest version of this
%% license is in:
%%
%%    http://www.latex-project.org/lppl.txt
%%
%% and version 1.2 or later is part of all distributions of LaTeX version
%% 1999/12/01 or later.
%%
%% In contrast to the original, we do not include the definitions for/using:
%% gridpicture, div[2], isMEMOIR[1], gridSetup[6][], subgridstyle{dotted}, labelfactor{}, gap{}, gridunitname{}, gridunit{}, gridlines{\thinlines}, subgridlines{\thinlines}, the {keyval} package, evenside margin, nor any definitions with 'color'.
%%
%% These are beyond  what is needed for the NAACL style.
%%
\newcommand\LenToUnit[1]{#1\@gobble}
\newcommand\AtPageUpperLeft[1]{%
  \begingroup
    \@tempdima=0pt\relax\@tempdimb=\ESO@yoffsetI\relax
    \put(\LenToUnit{\@tempdima},\LenToUnit{\@tempdimb}){#1}%
  \endgroup
}
\newcommand\AtPageLowerLeft[1]{\AtPageUpperLeft{%
  \put(0,\LenToUnit{-\paperheight}){#1}}}
\newcommand\AtPageCenter[1]{\AtPageUpperLeft{%
  \put(\LenToUnit{.5\paperwidth},\LenToUnit{-.5\paperheight}){#1}}}
\newcommand\AtPageLowerCenter[1]{\AtPageUpperLeft{%
  \put(\LenToUnit{.5\paperwidth},\LenToUnit{-\paperheight}){#1}}}%
\newcommand\AtPageLowishCenter[1]{\AtPageUpperLeft{%
  \put(\LenToUnit{.5\paperwidth},\LenToUnit{-.96\paperheight}){#1}}}
\newcommand\AtTextUpperLeft[1]{%
  \begingroup
    \setlength\@tempdima{1in}%
    \advance\@tempdima\oddsidemargin%
    \@tempdimb=\ESO@yoffsetI\relax\advance\@tempdimb-1in\relax%
    \advance\@tempdimb-\topmargin%
    \advance\@tempdimb-\headheight\advance\@tempdimb-\headsep%
    \put(\LenToUnit{\@tempdima},\LenToUnit{\@tempdimb}){#1}%
  \endgroup
}
\newcommand\AtTextLowerLeft[1]{\AtTextUpperLeft{%
  \put(0,\LenToUnit{-\textheight}){#1}}}
\newcommand\AtTextCenter[1]{\AtTextUpperLeft{%
  \put(\LenToUnit{.5\textwidth},\LenToUnit{-.5\textheight}){#1}}}
\newcommand{\ESO@HookI}{} \newcommand{\ESO@HookII}{}
\newcommand{\ESO@HookIII}{}
\newcommand{\AddToShipoutPicture}{%
  \@ifstar{\g@addto@macro\ESO@HookII}{\g@addto@macro\ESO@HookI}}
\newcommand{\ClearShipoutPicture}{\global\let\ESO@HookI\@empty}
\newcommand{\@ShipoutPicture}{%
  \bgroup
    \@tempswafalse%
    \ifx\ESO@HookI\@empty\else\@tempswatrue\fi%
    \ifx\ESO@HookII\@empty\else\@tempswatrue\fi%
    \ifx\ESO@HookIII\@empty\else\@tempswatrue\fi%
    \if@tempswa%
      \@tempdima=1in\@tempdimb=-\@tempdima%
      \advance\@tempdimb\ESO@yoffsetI%
      \unitlength=1pt%
      \global\setbox\@cclv\vbox{%
        \vbox{\let\protect\relax
          \pictur@(0,0)(\strip@pt\@tempdima,\strip@pt\@tempdimb)%
            \ESO@HookIII\ESO@HookI\ESO@HookII%
            \global\let\ESO@HookII\@empty%
          \endpicture}%
          \nointerlineskip%
        \box\@cclv}%
    \fi
  \egroup
}
\EveryShipoutTACL{\@ShipoutPicture}
\newif\ifESO@dvips\ESO@dvipsfalse
\newif\ifESO@grid\ESO@gridfalse
\newif\ifESO@texcoord\ESO@texcoordfalse
\newcommand*\ESO@griddelta{}\newcommand*\ESO@griddeltaY{}
\newcommand*\ESO@gridDelta{}\newcommand*\ESO@gridDeltaY{}
\newcommand*\ESO@yoffsetI{}\newcommand*\ESO@yoffsetII{}
\ifESO@texcoord
  \def\ESO@yoffsetI{0pt}\def\ESO@yoffsetII{-\paperheight}
  \edef\ESO@griddeltaY{-\ESO@griddelta}\edef\ESO@gridDeltaY{-\ESO@gridDelta}
\else
  \def\ESO@yoffsetI{\paperheight}\def\ESO@yoffsetII{0pt}
  \edef\ESO@griddeltaY{\ESO@griddelta}\edef\ESO@gridDeltaY{\ESO@gridDelta}
\fi


%% ----- Submitted version markup: Page numbers, ruler, and confidentiality.  Using ideas/code from cvpr.sty 2015. --MM

\font\naaclhv  = phvb at 8pt

%% Define vruler %%

%\makeatletter
\newbox\aclrulerbox
\newcount\aclrulercount
\newdimen\aclruleroffset
\newdimen\cv@lineheight
\newdimen\cv@boxheight
\newbox\cv@tmpbox
\newcount\cv@refno
\newcount\cv@tot
% NUMBER with left flushed zeros  \fillzeros[<WIDTH>]<NUMBER>
\newcount\cv@tmpc@ \newcount\cv@tmpc
\def\fillzeros[#1]#2{\cv@tmpc@=#2\relax\ifnum\cv@tmpc@<0\cv@tmpc@=-\cv@tmpc@\fi
\cv@tmpc=1 %
\loop\ifnum\cv@tmpc@<10 \else \divide\cv@tmpc@ by 10 \advance\cv@tmpc by 1 \fi
   \ifnum\cv@tmpc@=10\relax\cv@tmpc@=11\relax\fi \ifnum\cv@tmpc@>10 \repeat
\ifnum#2<0\advance\cv@tmpc1\relax-\fi
\loop\ifnum\cv@tmpc<#1\relax0\advance\cv@tmpc1\relax\fi \ifnum\cv@tmpc<#1 \repeat
\cv@tmpc@=#2\relax\ifnum\cv@tmpc@<0\cv@tmpc@=-\cv@tmpc@\fi \relax\the\cv@tmpc@}%
% \makevruler[<SCALE>][<INITIAL_COUNT>][<STEP>][<DIGITS>][<HEIGHT>]
\def\makevruler[#1][#2][#3][#4][#5]{\begingroup\offinterlineskip
\textheight=#5\vbadness=10000\vfuzz=120ex\overfullrule=0pt%
\global\setbox\aclrulerbox=\vbox to \textheight{%
{\parskip=0pt\hfuzz=150em\cv@boxheight=\textheight
\cv@lineheight=#1\global\aclrulercount=#2%
\cv@tot\cv@boxheight\divide\cv@tot\cv@lineheight\advance\cv@tot2%
\cv@refno1\vskip-\cv@lineheight\vskip1ex%
\loop\setbox\cv@tmpbox=\hbox to0cm{{\naaclhv\hfil\fillzeros[#4]\aclrulercount}}%
\ht\cv@tmpbox\cv@lineheight\dp\cv@tmpbox0pt\box\cv@tmpbox\break
\advance\cv@refno1\global\advance\aclrulercount#3\relax
\ifnum\cv@refno<\cv@tot\repeat}}\endgroup}%
%\makeatother

\def\confidentialtext{Confidential TACL submission. DO NOT DISTRIBUTE.}
\def\confidential{\confidentialtext}

%% Page numbering, Vruler and Confidentiality %%
% \makevruler[<SCALE>][<INITIAL_COUNT>][<STEP>][<DIGITS>][<HEIGHT>]

% SC/KG/WL - changed line numbering to gainsboro
\definecolor{gainsboro}{rgb}{0.8, 0.8, 0.8}
%\def\aclruler#1{\makevruler[14.17pt][#1][1][3][\textheight]\usebox{\aclrulerbox}}  %% old line
\def\aclruler#1{\textcolor{gainsboro}{\makevruler[14.17pt][#1][1][3][\textheight]\usebox{\aclrulerbox}}}

\def\leftoffset{-2.1cm} %original: -45pt
\def\rightoffset{17.5cm} %original: 500pt
\iftaclpubformat\else\pagenumbering{arabic}\fi
\AddToShipoutPicture{%
\iftaclpubformat\else
\AtPageLowishCenter{\thepage}
\aclruleroffset=\textheight
\advance\aclruleroffset4pt
  \AtTextUpperLeft{%
    \put(\LenToUnit{\leftoffset},\LenToUnit{-\aclruleroffset}){%left ruler
      \aclruler{\aclrulercount}}
    \put(\LenToUnit{\rightoffset},\LenToUnit{-\aclruleroffset}){%right ruler
      \aclruler{\aclrulercount}}
  }
  \AtTextUpperLeft{%confidential
    \put(0,\LenToUnit{1cm}){\parbox{\textwidth}{\centering\naaclhv\confidential}}
  }
\fi
}

%%%% ----- End settings for placing additional items into the submitted version --MM ----- %%%%

%%%% ----- Begin settings for both submitted and camera-ready version ----- %%%%

%% Title and Authors %%

\newcommand{\anonsubtext}{Anonymous TACL submission}
\newcommand\outauthor{
    \begin{tabular}[t]{c}
  \iftaclpubformat
       \bf\@author
  \else
    % Avoiding common accidental de-anonymization issue. --MM
     \bf \anonsubtext
  \fi
    \end{tabular}}


% Changing the expanded titlebox for submissions to 2.5 in (rather than 6.5cm)
% and moving it to the style sheet, rather than within the example tex file. --MM
\iftaclpubformat
\else
  \addtolength\titlebox{.25in}
\fi
% Mostly taken from deproc.
\def\maketitle{\par
 \begingroup
   \def\thefootnote{\fnsymbol{footnote}}
   \def\@makefnmark{\hbox to 0pt{$^{\@thefnmark}$\hss}}
   \twocolumn[\@maketitle] \@thanks
 \endgroup
 \setcounter{footnote}{0}
 \let\maketitle\relax \let\@maketitle\relax
 \gdef\@thanks{}\gdef\@author{}\gdef\@title{}\let\thanks\relax
 \iftaclcopyedit
  \doublespacing % after title and author, start doublespacing
 \fi
 }
\def\@maketitle{\vbox to \titlebox{\hsize\textwidth
 \linewidth\hsize \vskip 0.125in minus 0.125in \centering
 {\Large\bf \@title \par} \vskip 0.2in plus 1fil minus 0.1in
 {\def\and{\unskip\enspace{\rm and}\enspace}%
  \def\And{\end{tabular}\hss \egroup \hskip 1in plus 2fil
           \hbox to 0pt\bgroup\hss \begin{tabular}[t]{c}\bf}%
  \def\AND{\end{tabular}\hss\egroup \hfil\hfil\egroup
          \vskip 0.25in plus 1fil minus 0.125in
           \hbox to \linewidth\bgroup\large \hfil\hfil
             \hbox to 0pt\bgroup\hss \begin{tabular}[t]{c}\bf}
  \hbox to \linewidth\bgroup\large \hfil\hfil
    \hbox to 0pt\bgroup\hss
  \outauthor
   \hss\egroup
    \hfil\hfil\egroup}
  \vskip 0.3in plus 2fil minus 0.1in
}}

% margins for abstract
\renewenvironment{abstract}{\centerline{\large\bf
 Abstract}\vspace{0.5ex}\begin{quote} \small}{\par\end{quote}\vskip 1ex}

\RequirePackage{natbib}
% for citation commands in the .tex, authors can use:
% \citep, \citet, and \citeyearpar for compatibility with natbib, or
% \cite, \newcite, and \shortcite for compatibility with older ACL .sty files
\renewcommand\cite{\citep}  % to get "(Author Year)" with natbib
\newcommand\shortcite{\citeyearpar}% to get "(Year)" with natbib
\newcommand\newcite{\citet} % to get "Author (Year)" with natbib


% bibliography

\def\@up#1{\raise.2ex\hbox{#1}}

% Don't put a label in the bibliography at all.  Just use the unlabeled format
% instead.
\def\thebibliography#1{\vskip\parskip%
\vskip\baselineskip%
\def\baselinestretch{1}%
\ifx\@currsize\normalsize\@normalsize\else\@currsize\fi%
\vskip-\parskip%
\vskip-\baselineskip%
\section*{References\@mkboth
 {References}{References}}\list
 {}{\setlength{\labelwidth}{0pt}\setlength{\leftmargin}{\parindent}
 \setlength{\itemindent}{-\parindent}}
 \def\newblock{\hskip .11em plus .33em minus -.07em}
 \sloppy\clubpenalty4000\widowpenalty4000
 \sfcode`\.=1000\relax}
\let\endthebibliography=\endlist


% Allow for a bibliography of sources of attested examples
\def\thesourcebibliography#1{\vskip\parskip%
\vskip\baselineskip%
\def\baselinestretch{1}%
\ifx\@currsize\normalsize\@normalsize\else\@currsize\fi%
\vskip-\parskip%
\vskip-\baselineskip%
\section*{Sources of Attested Examples\@mkboth
 {Sources of Attested Examples}{Sources of Attested Examples}}\list
 {}{\setlength{\labelwidth}{0pt}\setlength{\leftmargin}{\parindent}
 \setlength{\itemindent}{-\parindent}}
 \def\newblock{\hskip .11em plus .33em minus -.07em}
 \sloppy\clubpenalty4000\widowpenalty4000
 \sfcode`\.=1000\relax}
\let\endthesourcebibliography=\endlist

% sections with less space
\def\section{\@startsection {section}{1}{\z@}{-2.0ex plus
    -0.5ex minus -.2ex}{1.5ex plus 0.3ex minus .2ex}{\large\bf\raggedright}}
\def\subsection{\@startsection{subsection}{2}{\z@}{-1.8ex plus
    -0.5ex minus -.2ex}{0.8ex plus .2ex}{\normalsize\bf\raggedright}}
%% changed by KO to - values to get teh initial parindent right
\def\subsubsection{\@startsection{subsubsection}{3}{\z@}{-1.5ex plus
   -0.5ex minus -.2ex}{0.5ex plus .2ex}{\normalsize\bf\raggedright}}
\def\paragraph{\@startsection{paragraph}{4}{\z@}{1.5ex plus
   0.5ex minus .2ex}{-1em}{\normalsize\bf}}
\def\subparagraph{\@startsection{subparagraph}{5}{\parindent}{1.5ex plus
   0.5ex minus .2ex}{-1em}{\normalsize\bf}}

% Footnotes
\footnotesep 6.65pt %
\skip\footins 9pt plus 4pt minus 2pt
\def\footnoterule{\kern-3pt \hrule width 5pc \kern 2.6pt }
\setcounter{footnote}{0}

% Lists and paragraphs
\parindent 1em
\topsep 4pt plus 1pt minus 2pt
\partopsep 1pt plus 0.5pt minus 0.5pt
\itemsep 2pt plus 1pt minus 0.5pt
\parsep 2pt plus 1pt minus 0.5pt

\leftmargin 2em \leftmargini\leftmargin \leftmarginii 2em
\leftmarginiii 1.5em \leftmarginiv 1.0em \leftmarginv .5em \leftmarginvi .5em
\labelwidth\leftmargini\advance\labelwidth-\labelsep \labelsep 5pt

\def\@listi{\leftmargin\leftmargini}
\def\@listii{\leftmargin\leftmarginii
   \labelwidth\leftmarginii\advance\labelwidth-\labelsep
   \topsep 2pt plus 1pt minus 0.5pt
   \parsep 1pt plus 0.5pt minus 0.5pt
   \itemsep \parsep}
\def\@listiii{\leftmargin\leftmarginiii
    \labelwidth\leftmarginiii\advance\labelwidth-\labelsep
    \topsep 1pt plus 0.5pt minus 0.5pt
    \parsep \z@ \partopsep 0.5pt plus 0pt minus 0.5pt
    \itemsep \topsep}
\def\@listiv{\leftmargin\leftmarginiv
     \labelwidth\leftmarginiv\advance\labelwidth-\labelsep}
\def\@listv{\leftmargin\leftmarginv
     \labelwidth\leftmarginv\advance\labelwidth-\labelsep}
\def\@listvi{\leftmargin\leftmarginvi
     \labelwidth\leftmarginvi\advance\labelwidth-\labelsep}

\abovedisplayskip 7pt plus2pt minus5pt%
\belowdisplayskip \abovedisplayskip
\abovedisplayshortskip  0pt plus3pt%
\belowdisplayshortskip  4pt plus3pt minus3pt%

% Less leading in most fonts (due to the narrow columns)
% The choices were between 1-pt and 1.5-pt leading
% \def\@normalsize{\@setsize\normalsize{11pt}\xpt\@xpt}
% \def\small{\@setsize\small{10pt}\ixpt\@ixpt}
% \def\footnotesize{\@setsize\footnotesize{10pt}\ixpt\@ixpt}
% \def\scriptsize{\@setsize\scriptsize{8pt}\viipt\@viipt}
% \def\tiny{\@setsize\tiny{7pt}\vipt\@vipt}
% \def\large{\@setsize\large{14pt}\xiipt\@xiipt}
\def\Large{\@setsize\Large{16pt}\xivpt\@xivpt}
% \def\LARGE{\@setsize\LARGE{20pt}\xviipt\@xviipt}
% \def\huge{\@setsize\huge{23pt}\xxpt\@xxpt}
% \def\Huge{\@setsize\Huge{28pt}\xxvpt\@xxvpt}


\let\@@makecaption\@makecaption
\renewcommand{\@makecaption}[1]{\@@makecaption{\small #1}}

\end{filecontents*}
\usepackage{times,latexsym}
\usepackage[T1]{fontenc}
\usepackage{amsmath,amssymb,amsthm,booktabs}
\usepackage{microtype}
\microtypesetup{expansion=false}
\usepackage[]{tacl2021v1}
\hypersetup{pdfauthor={},pdftitle={Lossy Calibration Archives for Exact Model Repair},pdfsubject={Anonymous TACL submission}}
\newtheorem{theorem}{Theorem}
\newtheorem{lemma}{Lemma}
\newcommand{\Qfix}{Q_{\mathrm{fix}}}
\title{Lossy Calibration Archives for Exact Model Repair}
\author{Anonymous submission}
\date{}
\begin{document}
\maketitle
\begin{abstract}
Low-bit language models use calibration text to choose a small set of allowed
weight values. Removing calibration records can change those choices. Recovering
the revised model usually requires the original text or a large archive of its
internal features. We show how a lossy feature archive can still recover the exact
weight choices. The archive stores intervals, and a verifier accepts a choice
only when every feature value in those intervals leads to it. Unresolved choices
trigger source replay. For a calibration procedure whose feature extractor stays
fixed, we prove exact model recovery and equality of the state needed for later
deletions. For this target, DistilGPT2 and GPT2 repairs match fresh reconstruction across
42.47 and 84.93 million calibrated weight codes. In a matched study,
compression saves 7.01--7.86\% of live state, or 1.30--2.24\% including the
checkpoint, while taking 1.44--2.66 times the optimized lossless request time. The result is a checked storage and source-access
tradeoff: exact discrete model decisions can survive the loss of exact feature
values.
\end{abstract}

\section{Introduction}
A language model can be made smaller by replacing each weight with one of a few
allowed values. Calibration text helps choose these values by showing which
weight errors matter for the model's internal computations. That text is often
treated as a one-time input. Yet a calibration collection can later change. A
record may be duplicated, corrupted, or removed from a deployment's approved
collection. The resulting maintenance task is concrete: rebuild the low-bit
weights for the retained records and keep enough state to process the next
removal.

A service can retain the text and extract features again, cache exact features,
or aggregate their contributions into calibration matrices. Each choice trades
storage against later source access. An aggregate alone may not identify the
contribution that must be removed.

We study a fourth option: retaining compressed intervals around the feature
values. At first this seems incompatible with exact reconstruction. If compression
loses a feature value, how can a later computation be exact? The key is that the
required output is discrete. Many possible feature values lead to the same
rounded weight. The archive need only preserve enough information to distinguish
the relevant rounding choices. It need not reconstruct every original feature
bit.

Our verifier recomputes retained-set decisions, including changed codes. A
numerical solver proposes the calculation; directed bounds check it for every
feature matrix described by the archive. When this check cannot resolve a choice,
the service can regenerate exact retained-source features.

We study maintenance of a specified calibration procedure. A feature extractor
from the original checkpoint is fixed before choosing calibration membership.
Each record then has the same features regardless of which other records remain.
This \emph{fixed-teacher calibration} supports source-local archives and exact
subtraction of source contributions. Ordinary sequential calibration instead
allows earlier quantized layers to alter later features. Our target also fixes the ridge, normalization, and checkpoint-only grids.
We measure its predictive cost against a matched sequential constructor. The task changes calibration
membership; it does not remove information from the checkpoint's pretraining.

We contribute a repair service using lossy feature enclosures, executable
certificates for its adaptive decisions, and proofs of exact model and successor
state recovery. Complete transactions compare cold reconstruction, two lossless
caches, and an exact-Gram hybrid, separating storage, source access, time, and
prediction quality.

Adaptive rounding already uses calibration geometry~\citep{gptq,quip}; verified
numerics and adaptive predicates already certify discrete decisions~\citep{rump,shewchuk}.
Our focus is their use when the calibration evidence itself is lossy and the
result must support another deletion.

\section{What must the archive remember?}
\label{sec:intuition}
Consider a service maintaining a four-bit model. Four-bit weights use sixteen
allowed values, called \emph{codes}. Calibration adjusts one weight at a time.
An earlier rounding error influences the next choice through coefficients
computed from the calibration features. Deletion changes those coefficients,
which can change later codes. Exact repair therefore requires agreement along
this entire adaptive sequence.

\paragraph{A changed code from uncertain features.}
A two-weight example illustrates the information requirement. Let the original
weights be $(0.4,0.4)$, with allowed values $\{-1,0,1\}$. Suppose two records have
feature vectors $(1,1)$ and $(-1,1)$. Their cross-products cancel. The first weight
rounds to zero, and the second receives no correction, so the output is $(0,0)$.
Now delete the second record. With a ridge term of one, the coefficient carrying
the first rounding error into the second coordinate is
\[
 c=\frac{x_1x_2}{1+x_2^2}=\frac12.
\]
The second adjusted value becomes $0.4+c(0.4-0)=0.6$, and its code becomes one.

Suppose the archive retains only $x_1,x_2\in[0.99,1.01]$. Direct interval bounds give
\[
 \frac{0.9801}{2.0201}\leq c\leq\frac{1.0201}{1.9801}.
\]
Every resulting adjusted value lies inside $[0.59,0.61]$, entirely above the
rounding boundary $0.5$. Thus the revised output $(0,1)$ is exact even though the
archive cannot identify either original feature value. This is an algebraic
illustration, not an experimental result. If the interval crossed $0.5$, the same
check would request more information rather than guess a code.

\paragraph{Three separate objects.}
The model contains deployed codes. The repair state identifies retained records
and stores evidence for later deletion. The execution history records work,
including replay and failure. Identical codes need not imply identical repair
states, and identical states need not imply identical costs. We evaluate these
objects separately.

\section{The reconstruction target}
\label{sec:target}
We now define the model that repair must recover. The specification fixes a
pretrained checkpoint $W$, a set of identified token records $\mathcal D_0$, a
teacher $T(W)$, and the tokenizer and record boundaries. It also fixes stage and
coordinate order, rounding grids, positive ridge parameters, and the positive
original normalization $N_0$. Deleting records changes membership, not these choices.
The experiments use the original checkpoint as teacher, $T(W)=W$.

\subsection{Source-local features}
A stage is one calibrated affine transformation. For stage $\ell$ and record $j$,
let
\[
 X_{\ell j}=F_\ell(T(W),j)\in\mathbb Q^{d_\ell\times t_j},
\]
where $d_\ell$ is the feature width and $t_j$ is the record's token count.
$F_\ell$ is a declared finite CPU evaluation. Each successful binary64 output is
interpreted exactly as a dyadic rational, an integer times a power of two. This
specifies the arithmetic object being reconstructed without replacing the
feature extractor by an ideal real-arithmetic network. Unsupported finite
operations cause refusal.

The evaluator processes each record independently. For retained membership
$R\subseteq\mathcal D_0$, concatenate its features in canonical source order to
form $X_\ell(R)$. Its regularized second-moment matrix is
\begin{equation}
 H_\ell(R)=\lambda_\ell I+
 \frac{X_\ell(R)X_\ell(R)^\top}{N_0},\qquad\lambda_\ell>0.
 \label{eq:metric}
\end{equation}
Here $I$ is the identity matrix. The positive ridge makes the matrix positive
definite. Fixing $N_0$ keeps the unnormalized ridge $\beta=\lambda_\ell N_0$
constant across deletions. If $T$ tokens remain, multiplying $H$ by $N_0/T$
gives their average Gram plus ridge $\lambda_\ell N_0/T$. Thus the declared
policy strengthens regularization relative to the retained-data average as
records disappear; changing that policy defines a different target.

\subsection{Adaptive weight decisions}
Suppress stage and output-row indices. Factor the rational matrix as
$H=L^\top\operatorname{diag}(\tau)L$, where $L$ is unit lower triangular and the
pivots $\tau_i$ are positive. This reverse LDL factorization determines the
coefficient $c_{ih}=L_{ih}$ carrying an earlier rounding error into coordinate
$i$. For a base-weight row $w$, define
\begin{align}
 v_i&=w_i+\sum_{h<i}c_{ih}(w_h-q_h),\label{eq:recurrence}\\
 q_i&=\operatorname{nearest}_{G_i}(v_i).\label{eq:rounding}
\end{align}
This recurrence uses the established quadratic rounding viewpoint~\citep{quip}.
For a proposed real weight row
$y$, the distortion is $(w-y)H(w-y)^\top$. With $y_h=q_h$ for $h<i$, minimize
over the remaining real coordinates. The minimizing coordinate $i$ is $v_i$;
then round it to $G_i$. The suffix equations in Section~\ref{sec:coefficient}
give this identity. An exact midpoint selects the lower code, and endpoint
codes saturate. Collecting all stage outputs defines $\Qfix(R)$; this greedy
procedure need not minimize distortion jointly over all discrete codes.
The ridge term penalizes weight movement.

For each experimental output row, the grid is
$s_{\ell,r}\{-8,\ldots,7\}$. Its scale is the smallest positive checkpoint-only dyadic
scale with at most 24 significant binary digits that covers the base weights
and gives finite exactly representable codes and midpoints. The scale is shared
across that row's coordinates. Archive precision and weight precision are
separate: the archive stores feature evidence, while these sixteen values define
the deployed four-bit weights.

Earlier decisions enter later values through~\eqref{eq:recurrence}. A certificate
must therefore establish the accepted prefix before using it to certify the next
code. The fixed teacher keeps later defining inputs independent of repaired weights.
Changing any fixed target component creates a different reconstruction problem.

\section{Certifying decisions from lossy evidence}
\label{sec:method}
Preparation evaluates each source through the frozen teacher and encodes a
canonical descriptor $D_{\ell j}$. Decoding returns entrywise bounds satisfying
\begin{equation}
 \underline X_{\ell j}\leq X_{\ell j}\leq\overline X_{\ell j}.
 \label{eq:containment}
\end{equation}
Preparation checks every actual feature value against its bounds. Metadata binds
the source and token identities, stage, target, codec, and a commitment to the
original feature words.

The codec groups 256 token-major values into blocks. Each block stores a shared
power-of-two step and signed 48-bit cell indices. Two-bit tags distinguish closed
cells, exact grid values, and exact binary64 escapes near overflow. Endpoints are
exactly representable dyadics. Escapes and metadata count toward storage, so
48-bit indices do not imply a uniform 25\% archive saving.

\subsection{From coefficient intervals to codes}
A request first validates its predecessor, deletion IDs, retained membership,
and target bindings. It keeps each surviving descriptor unchanged and processes
stages in order. A numerical solver proposes coefficients. A separate directed
verifier bounds their errors over every feature realization in the decoded box.
Using the accepted prefix in~\eqref{eq:recurrence}, it obtains an enclosure
$[\underline v_i,\overline v_i]$ for the next adjusted value.

For a grid $g_0<\cdots<g_{k-1}$, let $m_a=(g_a+g_{a+1})/2$. The decision cell
for an interior code $g_a$ is $(m_{a-1},m_a]$. Endpoint cells extend to infinity.
Acceptance requires the entire value enclosure to lie in one cell. The strict
lower boundary and inclusive upper boundary implement the lower-code tie rule.
A proposal close to a grid value is insufficient if its verified enclosure
crosses a cell boundary.

If a box cannot certify the stage within its admitted work, the service may
regenerate exact retained features and use a point certificate. It verifies that
replayed words and containment match preparation. Access accounting includes
every traversed ancestor stage. Temporary exact features do not replace the
compressed successor descriptors. Exhausting a CPU, memory, refinement, or replay
allowance ends the attempt without committing a complete replacement.

\subsection{An executable coefficient bound}
\label{sec:coefficient}
We give the link between a numerical solve and the adaptive coefficients.
For one realized feature matrix $X$, set $K=XX^\top$, $\beta=\lambda N_0>0$, and
$A=\beta I+K=N_0H$. Let $S_i=\{i,\ldots,d\}$ be the suffix beginning at
coordinate $i$, and let $e_1$ be its first coordinate vector. Then
\begin{equation}
 c_{ih}=K_{h,S_i}A_{S_i,S_i}^{-1}e_1,\qquad h<i.
 \label{eq:suffix}
\end{equation}
To prove the identity, factor $A=L^\top D L$. Scaling $H$ changes the diagonal
factor, not $L$. With $S=S_i$, the suffix and cross blocks are
\begin{align*}
 A_{S,S}&=L_{S,S}^\top D_{S,S}L_{S,S},\\
 A_{h,S}&=L_{S,h}^\top D_{S,S}L_{S,S}.
\end{align*}
Their product with $A_{S,S}^{-1}e_1$ is
$L_{S,h}^\top L_{S,S}^{-\top}e_1=L_{ih}$: the first column of the unit upper
triangular inverse is $e_1$. Since $h<i$, $A_{h,S}=K_{h,S}$, proving~\eqref{eq:suffix}.

Let $\widehat z_i$ be a numerical proposal, held fixed throughout its check.
Suppose the verifier establishes
$\|e_1-A_{S,S}\widehat z_i\|_2\leq\eta_i$, where $\|\cdot\|_2$ is the Euclidean
norm. The coefficient error obeys
\begin{equation}
 \left|c_{ih}-K_{h,S}\widehat z_i\right|
 \leq\eta_i\min\left\{
 \frac{\|K_{h,S}\|_2}{\beta},
 \frac{\sqrt{K_{hh}}}{2\sqrt\beta}\right\}.
 \label{eq:radius}
\end{equation}
For the first bound, write the error as
$K_{h,S}A_{S,S}^{-1}(e_1-A_{S,S}\widehat z_i)$ and use
$A_{S,S}\succeq\beta I$. For the second, use $K_{h,S}=X_hX_S^\top$.
Each singular value $s$ of $X_S$ contributes $s/(\beta+s^2)$ to the norm of
$X_S^\top(\beta I+X_SX_S^\top)^{-1}$. This is at most $1/(2\sqrt\beta)$, while
$\|X_h\|_2=\sqrt{K_{hh}}$. Multiplication by the residual bound proves the claim.

For uncertain features, directed arithmetic encloses the dot product
$K_{h,S}\widehat z_i$, the residual norm, the Gram-row norm, and its diagonal
uniformly over the feature box. Combining this dot-product enclosure with
\eqref{eq:radius} bounds the coefficient. The argument applies to each realized
feature Gram, which is positive semidefinite. It does not assign that property
to arbitrary matrices inside an entrywise Gram enclosure. The coefficient
intervals then feed the decision-cell check above.

\subsection{How much uncertainty is tolerable?}
\label{sec:margin}
The decision check gives a useful criterion for archive precision. Suppose the
coefficient intervals for the accepted prefix have centers $\bar c_{ih}$ and
radii $\delta_{ih}\geq0$. Define the adjusted center and propagated radius as
\begin{align}
 u_i&=w_i+\sum_{h<i}\bar c_{ih}(w_h-q_h),\\
 \rho_i&=\sum_{h<i}\delta_{ih}|w_h-q_h|.
 \label{eq:propagation}
\end{align}
For every admitted coefficient realization,
$|v_i-u_i|\leq\rho_i$ by the triangle inequality. If a proposed code has cell
$(a_i,b_i]$, the sufficient acceptance conditions are
\begin{equation}
 u_i-\rho_i>a_i,\qquad u_i+\rho_i\leq b_i.
 \label{eq:margin}
\end{equation}
At endpoint codes the corresponding infinite boundary imposes no restriction.
These expressions are exact arithmetic identities; an executable evaluation
must enclose the centers, products, sums, and comparisons outwardly as well.

\paragraph{From archive width to coefficient error.}
Let $\bar X$ be the box center and $E$ its entrywise radius. Set
$\epsilon=\|E\|_F$ and $\kappa=\|\bar X\|_F$, where $\|\cdot\|_F$ is the
square root of the sum of squared entries. Let $\bar K=\bar X\bar X^\top$ and $\bar A=\beta I+\bar K$. Define
\[
 \gamma=2\kappa\epsilon+\epsilon^2,\qquad
 D_{ih}=\frac{\gamma}{\beta}
       +\frac{\|\bar K_{h,S_i}\|_2\gamma}{\beta^2}.
\]
Every feature matrix in the box satisfies $|c_{ih}-\bar c_{ih}|\leq D_{ih}$,
where $\bar c_{ih}$ is the exact coefficient at $\bar X$. To see this, write
$X=\bar X+\Delta$ with $\|\Delta\|_2\leq\epsilon$, giving
$\|K-\bar K\|_2\leq\gamma$. Both regularized suffix inverses have norm at most
$1/\beta$. The resolvent identity bounds their difference by $\gamma/\beta^2$;
expanding $K_{h,S}A_{S,S}^{-1}-\bar K_{h,S}\bar A_{S,S}^{-1}$ proves the bound.

Along the center's accepted prefix, inserting $D_{ih}$ for $\delta_{ih}$
in~\eqref{eq:propagation} gives a sufficient archive-width condition via
\eqref{eq:margin}. For block cell widths $\Delta_b$, $\epsilon^2$ is the sum of
$(\Delta_b/2)^2$ over non-exact entries; exact escapes contribute zero.
Smaller cells reduce this bound, but small ridge and large earlier residuals
amplify it. This conservative analytical bound is not a prediction of the
bounded implementation's success. If two realizations give different codes,
no uniform certificate exists. Even when all agree, interval dependence can
cause refusal. An exact midpoint has zero two-sided margin; refinement alone
need not resolve it.

\subsection{Choosing the numerical representation}
The same coefficient has the token-space form
\begin{equation}
 c_{ih}=X_h(\beta I+X_S^\top X_S)^{-1}X_i^\top.
 \label{eq:dual}
\end{equation}
This follows by multiplying
$X_S^\top(\beta I+X_SX_S^\top)=(\beta I+X_S^\top X_S)X_S^\top$ by the respective
inverses. The identities allow a solver to work in either space. For
$B_i=\beta I+X_S^\top X_S$ and a fixed proposal $\widehat t_i$, a verified bound
$\|X_i^\top-B_i\widehat t_i\|_2\leq\xi_i$ gives
\[
 |c_{ih}-X_h\widehat t_i|\leq\|X_h\|_2\xi_i/\beta.
\]
Indeed, $B_i\succeq\beta I$, so the inverse has norm at most $1/\beta$.
Outward enclosures make this bound uniform over a feature box. The token route
can propagate it through accumulated weighted features instead of materializing
every scalar coefficient.

For width $d$, retained tokens $T$, and output rows $r$, admission uses
structural work terms $dT^2$ for token-space coefficients and $d^2T+d^3$ for
feature-space coefficients, with $rd^2$ row work for the latter. These are
budget proxies, not bounds on the full verification cost. The service admits its arrays and bounded schedule
before feature access. Its original dispatcher prefers token space when
$4T\leq d$, and streamed feature space otherwise. The final matched study fixes
the token route for all point-valued solves. These are declared engineering
policies. Actual runtime also depends on interval width, decoding, repeated
verification, and replay. Compatible controls receive the same kernels and
route choices.

\section{Correct models and reusable successor states}
\label{sec:correctness}
\begin{center}
\fbox{\parbox{0.92\columnwidth}{\small
\textbf{Repair procedure.} Prepare each source by evaluating its features,
encoding a descriptor, and checking containment. On deletion:\\
1. Validate the target, predecessor, IDs, and work allowance.\\
2. Keep surviving descriptors; propose and verify each stage's codes.\\
3. If unresolved and permitted, replay retained sources, check their words,
and certify the point input. Otherwise stop without a completed successor.\\
4. Verify all stage extents, serialize state and trust records, and seal completion.
Publish a successor only after these checks.}}
\end{center}
The first guarantee concerns the model produced by a successful request.
\begin{theorem}[Exact code recovery]
\label{thm:codes}
Assume that features come from the frozen source-local evaluator and that trusted
descriptors satisfy~\eqref{eq:containment}. Each accepted certificate must be
valid for every feature realization in its supplied box, conditional on its
accepted prefix. Any fallback must reproduce the same finite features and
certify~\eqref{eq:recurrence}--\eqref{eq:rounding}. If the service commits every
required stage, the resulting model is $\Qfix(R)$.
\end{theorem}
\begin{proof}
Source locality makes each surviving feature matrix equal to its retained-only
reconstruction. It lies inside the prepared box. A valid first-coordinate
certificate therefore selects the target code. Suppose the earlier accepted
codes equal the target prefix. The next certificate's premise is then satisfied,
so its selected code also equals the target. Induction establishes each row.
A successful fallback uses the same features and target, giving the same
conclusion. The fixed teacher makes this argument apply to every stage without
changing later defining inputs. The complete extent check covers the entire
committed model.
\end{proof}

The state must also support the next request. Fix the codec, preparer,
representation rule, source ordering, provenance policy, and serialization.
Let $E(R)$ contain the retained canonical descriptors or the canonical exact
pooled Grams required by that representation. Define
\begin{equation}
 \begin{aligned}
 \mathcal S(R)=\operatorname{encode}\bigl(&\text{target},\Qfix(R),R,\\
                         &E(R),\text{provenance}\bigr).
 \end{aligned}
 \label{eq:state}
\end{equation}
Provenance here means retained-record identity and target/codec bindings.
Execution logs are external. \emph{Trust sidecars} are separate verification
records binding exact Gram payloads to source commitments. The canonical state
contains the fields in~\eqref{eq:state}; required sidecars are checked separately
and included in live bytes. Retrieval tokens or indexes and any source-access
obligations are also counted.

\begin{lemma}[Exact source subtraction]
\label{lem:subtract}
For fixed source-local features, $K(R)=\sum_{j\in R}X_jX_j^\top$.
Exact subtraction of trusted deleted-source contributions, with original
normalization unchanged, equals fresh retained-source accumulation.
\end{lemma}
\begin{proof}
Each source contribution is unchanged by other records' deletion. Finite sums
over disjoint source sets are additive in rational arithmetic. Subtracting the
removed summands leaves the retained sum. Canonical integer scaling and
serialization remove dependence on the order of additions and deletions.
\end{proof}

\begin{theorem}[Canonical successor equality]
\label{thm:state}
Under Theorem~\ref{thm:codes}, suppose descriptors are deterministic and retained
unchanged, Gram updates satisfy Lemma~\ref{lem:subtract}, and compared executions
use the same policies in~\eqref{eq:state}. A successful repair returns
$\mathcal S(R)$. Its bytes equal those of any successful certified fresh
reconstruction in that representation. The equality holds after every step of
a sequence of successful valid deletions.
\end{theorem}
\begin{proof}
The model agrees by Theorem~\ref{thm:codes}. Retained descriptors agree because
their source-local inputs and encoder agree. Pooled Grams agree by exact
subtraction and canonicalization. Membership, target, provenance, and encoding
policies are identical, so each serialized component agrees. Apply the same
argument at each successor.
\end{proof}

Equality is representation-specific and excludes execution histories.
The trusted base includes the finite evaluator, outward arithmetic, runtime and
immutable input provenance. Hash commitments detect changed bytes, not semantic
containment or a hostile host. Positive ridge ensures invertibility, not completion.

\section{Evaluation design}
\label{sec:design}
We evaluate three questions separately: whether a complete repair is exact,
what resources it consumes, and how the fixed-teacher target affects predictions.
The service studies use DistilGPT2, with approximately 82 million parameters, and
GPT2, with 124 million (\citealp{distilgpt2}; \citealp{gpt2card}). They calibrate 24 and 48 affine
stages, respectively. The packed outputs contain 42,467,328 and 84,934,656 codes;
unchanged checkpoint components remain necessary for inference. Calibration
records are English C4 web text~\citep{c4}, each with 128 tokens.

\subsection{Comparators and complete accounting}
The \emph{cached lossless} arm stores exact feature values. The \emph{modern
lossless} arm selects the smallest complete descriptor among raw, zlib6,
byte-shuffled zlib6, Zstandard3, and byte-shuffled Zstandard3 encodings, with
fixed-order ties. Every candidate encoder is charged. \emph{Cold reconstruction}
extracts retained features from source records. The \emph{hybrid exact-Gram} arm
stores exact Grams at eighteen width-768 stages and lossless features at six
width-3072 stages in DistilGPT2. Our \emph{compressed} arm stores the 48-bit interval
descriptors. All arms target the same codes and are compared with a fresh
retained-input reconstruction in their respective state representation.
The cold arm deliberately writes the lossless successor cache, explaining its
identical stored size. It is not a minimal-storage source-only service; such a
baseline, retaining only codes, tokens or references, and bindings, is unmeasured.
It would require retained sources on every repair. Lossless caches need neither
retained nor deleted sources; compression may need retained sources on refusal;
the pooled-Gram hybrid may need deleted sources for subtraction.

An oracle traverses retained tokens without reading the candidate predecessor.
The audits compare actual model and state payloads; the later studies also compare
required trust artifacts. The oracle shares the numerical solver. This checks
complete reconstruction and source independence, while leaving implementation
errors shared by that solver outside the comparison.

The baseline request clock includes predecessor reads and validation, attempted
certificates, accumulation, extraction, fallback, output writes, and cleanup.
It excludes the caller's context loading and controller overhead. Optimization
and matched tables instead use complete host time, including startup and
admission. Comparisons are within each table; absolute times across studies
are not paired speed measurements. Nested source extraction is included in its
stage clock and is never added a second time.

Preparation must be included when comparing deployment lifetimes. If $P_a$ is
preparation for representation $a$ and $C_{a,k}$ is its complete cost on actual
successor transition $k$, then
\begin{equation}
 L_a(m)=P_a+\sum_{k=1}^{m}C_{a,k}.
 \label{eq:lifetime}
\end{equation}
Alternative deletions from an unchanged root do not form such a trajectory.
Shared preparation has an explicit attribution convention rather than an
implicit zero cost.

Live storage includes codes, descriptors or Grams, targets, tokens or retrieval
indexes, provenance, and required trust sidecars. We also report a total that
adds the common checkpoint. A model embedded in the state is counted once.
Runtime and codec installations, and physical audit archives, are separate.
Source access counts source--stage record traversals, including ancestors.
One record replayed through twenty-four stages costs twenty-four accesses.
This counts feature traversal, not distinct disk reads or network requests.

For a service with $n$ independently stored repair states and one shared extra
codec installation of $I$ bytes, compression saves total space only if
$\sum_{j=1}^n(B_{L,j}-B_{C,j})>I$, where $B_L,B_C$ count each complete state.
In Table~\ref{tab:service}, the modern-lossless differences are 33.959 and
16.997\,MB. Charging the entire 23.871\,MB codec installation as additional gives
10.089\,MB saved at twelve records and 6.873\,MB extra at six. Two separate
six-record states sharing that installation would save 10.124\,MB. This is a
sharing calculation, not a measured deployment. Repair frequency and retention
duration determine whether saved byte-time justifies extra compute; no service
cost model or deletion frequency was measured.

\subsection{Four studies and their selection history}
\label{sec:studies}
The first complete service study uses thirteen previously unexposed training
records. It deletes one and then six more, giving $13\to12\to6$ records while
keeping normalization 1664 and ridge $1/100$. Records are sorted by ID; the
first is deleted, followed by the next six. It compares all five arms. The second study holds
that root and target fixed and evaluates reference and optimized versions of the
compressed and modern lossless arms. Both optimized arms receive the compatible
token-neighbor kernel. Compression adds a native decoder; modern lossless adds
bounded reuse of decoded descriptors. Each package is evaluated as a whole.

The third study matches both checkpoints on four exposed records, deletion
$4\to3\to2$, normalization 512, and ridge $1/100$. Every point-valued solve uses the existing
token-space coefficient route, fixed before execution for all four cells and both fresh oracles.
Box and Gram policies remain unchanged. Each request runs in a fresh process on
the same host with one CPU and 16\,GiB memory. The grid-storage guard permits
1,327,104 scalar entries. Sealed preparation artifacts are reused and their costs retained.
All timing cells in these studies are single observations. Execution order is
fixed, with endpoint reversal in the optimization screen; OS cache is uncontrolled.

The fourth study evaluates the earlier 1664-normalized DistilGPT2 target on
untouched articles. It is distinct from the smaller matched targets. Teacher
and precision choices preceded it: a nearest-rounded teacher had failed a pooled
quality guard with a 5.92\% perplexity increase; the original-checkpoint teacher
was initially selected using eight exposed articles. A 40-bit archive refused
certification, followed by a separately registered 48-bit study. The initial
matched GPT2 request also refused after both preparations completed. A separate
stage diagnostic motivated the prospective point-token policy above. These
choices make the later studies scoped tests of fixed configurations.

\section{Complete repair results}
\label{sec:results}
\subsection{Storage and source access at a fixed target}
All ten arm--endpoint comparisons in the disjoint-root service study match fresh
reconstruction for the complete model and successor state. Table~\ref{tab:service}
shows their resource choices. Compression reduces live storage relative to both
lossless caches and needs no source traversal on this trajectory. Relative to
modern lossless, live savings are 9.24\% and 8.70\%; adding the checkpoint reduces
these to 4.71\% and 3.10\%. The codec installation adds 23,870,575 bytes separately.

\begin{table*}[t]
\centering\small
\begin{tabular}{rlrrrr}
\toprule
Records & Service arm & Request (s) & Live (MB) & With checkpoint (MB) & Accesses \\
\midrule
12 & Cached lossless & 4,419.007 & 368.003 & 720.828 & 0 \\
 & Modern lossless & 4,579.496 & 367.456 & 720.282 & 0 \\
 & Cold reconstruction & 5,973.008 & 368.003 & 720.828 & 288 \\
 & Hybrid exact Gram & 4,921.935 & 318.918 & 671.743 & 24 \\
 & Compressed & 5,206.569 & 333.497 & 686.322 & 0 \\
\midrule
6 & Cached lossless & 922.484 & 195.615 & 548.441 & 0 \\
 & Modern lossless & 897.841 & 195.330 & 548.155 & 0 \\
 & Cold reconstruction & 1,619.473 & 195.615 & 548.441 & 144 \\
 & Hybrid exact Gram & 1,694.311 & 220.247 & 573.072 & 144 \\
 & Compressed & 1,351.098 & 178.333 & 531.158 & 0 \\
\bottomrule
\end{tabular}
\caption{Complete DistilGPT2 service on $13\to12\to6$ records. MB are decimal;
live storage includes trust sidecars. Request time excludes caller context loading
and controller overhead. Accesses count source--stage traversals. The checkpoint
column excludes runtime and codec installation. Each row is one observation;
every model/state comparison passes.}
\label{tab:service}
\end{table*}

The hybrid demonstrates why the access contract matters. At twelve records, it
is both smaller and faster than compression, while reading source records. At
six, compression is smaller and faster than the hybrid and avoids its 144
accesses. Both lossless caches are faster than compression at both endpoints.
The compressed requests use bounded numerical fallback on 68 and 2 rows,
respectively, but no source replay. Numerical refinement and source access are
therefore distinct costs.

\subsection{Optimizing both sides of the comparison}
All eight requests in Table~\ref{tab:optimization} match fresh model, state, and
trust artifacts, with no source replay. The compressed package reduces observed
request times by 11.53\% and 36.17\%; the modern package reduces them by 0.69\%
and 19.93\%. The first optimized endpoint is nearly tied. The second favors
modern lossless. Their stored-state sizes are unchanged from
Table~\ref{tab:service}.

\begin{table}[t]
\centering\small
\begin{tabular}{@{}rlrr@{}}
\toprule
Records & Archive & Ref. (s) & Opt. (s) \\
\midrule
12 & Compressed & 4,444.116 & 3,931.646 \\
 & Modern & 3,995.241 & 3,967.523 \\
6 & Compressed & 1,328.624 & 848.119 \\
 & Modern & 939.553 & 752.336 \\
\bottomrule
\end{tabular}
\caption{Complete host times for both optimization packages on the same
DistilGPT2 trajectory. Modern denotes the modern lossless cache. Each timing cell is one request.}
\label{tab:optimization}
\end{table}

Including attributed shared preparation and both requests, the totals are
10,995.918 seconds for compression and 10,287.226 for modern lossless, a 6.89\%
observed lifetime penalty. The bundled changes do not isolate a decoder-only effect.

\subsection{Matched repair across two checkpoint depths}
All sixteen complete model, state, and trust comparisons pass for the matched
DistilGPT2/GPT2 study. Both audits cover every calibrated stage and code, including
the second successor. Table~\ref{tab:matched} reports all request times. The
optimization improves each representation relative to its reference, but
optimized compression takes longer than optimized modern lossless at all four
checkpoint/endpoint settings.

\begin{table*}[t]
\centering\small
\begin{tabular}{llrrrrr}
\toprule
Checkpoint & Records & Ref. comp. & Opt. comp. & Ref. lossless & Opt. lossless & Opt. C/L \\
\midrule
DistilGPT2 & 3 & 589.239 & 411.279 & 369.644 & 286.419 & 1.4359 \\
 & 2 & 436.691 & 325.131 & 219.802 & 179.706 & 1.8092 \\
GPT2 & 3 & 2,027.290 & 1,544.049 & 733.046 & 580.147 & 2.6615 \\
 & 2 & 838.153 & 615.039 & 406.506 & 326.341 & 1.8847 \\
\bottomrule
\end{tabular}
\caption{Complete host times (seconds) for the matched $4\to3\to2$ point-token
study. Lossless denotes the modern cache. C/L divides optimized compressed time
by optimized lossless time. Each cell is one observation, with complete equality
against its same-checkpoint fresh reconstruction.}
\label{tab:matched}
\end{table*}

\begin{table}[t]
\centering\small
\begin{tabular}{@{}lrrrr@{}}
\toprule
 & \multicolumn{2}{c}{Live (\%)} & \multicolumn{2}{c}{Total (\%)} \\
Records & 3 & 2 & 3 & 2 \\
\midrule
DistilGPT2 & 7.78 & 7.01 & 1.84 & 1.30 \\
GPT2 & 7.86 & 7.09 & 2.24 & 1.61 \\
\bottomrule
\end{tabular}
\caption{Matched compressed storage savings relative to modern lossless. Total
adds the shared checkpoint to each live-state size; it excludes runtime installs.}
\label{tab:storage}
\end{table}

Storage savings persist across both depths (Table~\ref{tab:storage}). Source
avoidance is less uniform. At three GPT2 records, both compressed cells replay
144 source--stage record pairs after refusal at the final-block MLP
up-projection. Optimized compressed stage work takes 1,458.767 seconds, including
695.104 seconds of source extraction, versus 476.578 seconds for modern lossless.
Six of eight compressed requests use no source replay; all eight lossless requests use none.

At two DistilGPT2 records, compression remains slower without replay: stage
work takes 267.560 seconds versus 114.651 for modern lossless. Thus cheaper
validation does not remove the certification cost. Both checkpoints share the
GPT family; complete code counts measure output coverage, not independent trials.

\subsection{Predictive cost of the fixed teacher}
\label{sec:quality}
Quality is evaluated for the earlier DistilGPT2 target with normalization 1664.
We select 320 articles without observing losses from 2,778 eligible records in
a saved frame of 3,801 complete records. The entire inspected frame, sixty
historical exclusions, and the separate calibration frame are protected against
reuse. Each article supplies 127 predictions from 128 tokens. Six model roles
share these articles: the fixed-teacher and sequential models at each retained
endpoint, the original checkpoint, and a nearest-grid control.

The matched sequential constructor uses the same retained records, normalization,
ridge, and grids, but obtains later features from its quantized prefix. It
isolates the teacher choice within our target family; it is not a comparison
with an off-the-shelf GPTQ implementation. A fixed-seed paired article bootstrap
uses 20,000 draws shared across endpoints. For article mean negative log losses
$\ell_{f,j}$ and $\ell_{s,j}$, the pooled ratio is
$\exp(\frac{1}{320}\sum_j(\ell_{f,j}-\ell_{s,j}))$, since articles have equal
prediction counts. The maximum is $\max_j\exp(\ell_{f,j}-\ell_{s,j})$.
Resampling preserves whole-article pairs; upper bounds use the empirical
97.5th percentile with the next higher order statistic.
The frozen guards require each pooled perplexity ratio and its
one-sided 97.5\% upper bound to be at most 1.05, and each observed maximum article
ratio to be at most 1.20. The two upper bounds have nominal familywise level
95\%. This is an approximate bootstrap over the admitted frame; the observed
maximum describes those articles.

\begin{table}[t]
\centering\small
\begin{tabular}{@{}rrrr@{}}
\toprule
Records & Pooled ratio & Upper bound & Maximum \\
\midrule
12 & 1.005421 & 1.010342 & 1.129140 \\
6 & 1.007149 & 1.011431 & 1.126961 \\
\bottomrule
\end{tabular}
\caption{Untouched quality: fixed-teacher perplexity divided by matched
sequential perplexity on 320 articles. The upper bound is one-sided 97.5\%.
All prespecified guards pass.}
\label{tab:quality}
\end{table}

The ratios in Table~\ref{tab:quality} show a small predictive cost. Candidate
perplexities are 60.1053 and 60.3088, compared with sequential values 59.7812 and
59.8807. The paired 95\% percentile intervals for the ratios are
[1.000455, 1.010342] and [1.002870, 1.011431], both above one. Original-checkpoint
and nearest-grid control perplexities are 54.9643 and 74.6245. These measurements
support the specified tolerance for this target and article frame. The
512-normalized targets in the matched study have no quality evaluation.

All twenty-four CPU/CUDA reference comparisons pass the fixed $10^{-8}$ mean
negative-log-likelihood barrier, with maximum difference $1.155\times10^{-14}$.
The terminal audit checks 1,920 GPU and twenty-four CPU losses, saved sequential
codes, original parameters, and recomputed statistics. This evaluates identical
saved models across likelihood implementations; calibration itself uses the
finite CPU evaluator defined in Section~\ref{sec:target}.

\section{Relation to prior work}
\label{sec:related}
\paragraph{Calibration and retained information.}
GPTQ~\citep{gptq} uses asymmetric min--max grids and propagates quantized
block outputs. Our fixed teacher and base-only grids define a different target.
QSS~\citep{qss} uses a frozen schema and additive prediction corrections;
removing schema-training data requires rebuilding the affected instance.
Our uncertain Grams instead determine adaptive weight codes.

QuIP~\citep{quip} develops LDL-based adaptive rounding and analyzes its
quadratic error. \citet{babai} relate reverse-order GPTQ to Babai's nearest-plane
algorithm, with a distortion bound under no clipping. Those results concern
rounding with known calibration geometry. We certify a fixed rounding target
when deletion changes that geometry and only feature enclosures remain.

\paragraph{Quantization and exact deletion.}
\citet{ginart} use quantized $k$-means and stability tests, with retraining
when deletion changes relevant centroids. Exact-Fun~\citep{exactfun} corrects
deleted gradients and resumes affected quantized federated training.
Our certificate covers changed retained-set codes over feature enclosures;
these different learning tasks are not runtime baselines.

\paragraph{Compression and verification.}
MEDU~\citep{medu} stores compressed distributed-training updates and bounds
deviation from retraining. Our output contract is equality of discrete codes.
QEBVerif~\citep{qebverif} bounds inference-output differences between a network
and its quantized counterpart using reachability analysis and MILP.
ExecCert~\citep{execcert} uses residual certificates and margins for discrete
prediction agreement. Our chain verifies adaptive weight codes and canonical
repair states. Interval bounds and proposal--verification separation are
established tools, rather than separate novelty claims.

\section{Reproducibility and practical use}
\label{sec:repro}
\paragraph{Implementation and data.}
The CPU runtime uses CPython 3.12.14, NumPy 2.3.5, OpenBLAS 0.3.30,
gmpy2 2.3.2, tokenizers 0.23.2, GCC 12.2, and Zstandard 1.5.7 through python-zstandard 0.25.0
on x86-64 Linux. Quality inference uses an NVIDIA A100 80GB; bootstrap statistics
use NumPy 2.2.6. Model citations pin checkpoint revisions.
The C4 revision is \texttt{1588ec454efa}. The thirteen-record root comes from
complete records in the first 2\,MiB compressed prefix of English training
shard 2. Quality uses validation shard 1 with separate 8\,MiB transfer and
decoded-frame bounds.
Records use the first 128 tokens without beginning/end markers, with positions
reset per source. Selection excludes exact IDs, text/normalized-text/URL hashes,
and token prefixes of lengths 16, 32, 64, and 128. Sampling is within these bounded
frames, not uniform across C4. Earlier development included WikiText-2~\citep{wikitext}.

Independent rational checks cover suffix/dual identities, conditional quadratic
minimization and cell boundaries, plus 39,366 realizations of 486 two-by-two
feature boxes for the archive-width bound. These are generated algebra checks,
separate from model experiments.

\paragraph{Executable arithmetic.}
Interval products enclose all endpoint products; sums accumulate outward at each
step using adjacent binary64 numbers. Exact zero shortcuts are checked explicitly.
Squared norm bounds accumulate upper magnitudes. A square-root proposal is
accepted only after its squared lower bound covers the normalized argument;
otherwise it is advanced upward, at most eight times. Overflow, nonfinite
bounds, and failed runtime checks cause refusal. Gradual underflow is required.
The native compiler disables fast math and fused contraction. Approximate
inverses and BLAS solves propose values; directed checks authorize codes.
Point refinement is capped at sixteen coordinates, with exact fallback only
within its separate rank and coordinate bounds. A midpoint uses the lower code
only when its exact value or full enclosure establishes that decision cell.

Feature evaluation retains coordinate-order binary64 products and sums, without
BLAS reduction or fused multiply-add. Nonlinear primitives use directed MPFR
bounds whose endpoints must round to the same binary64 word; an unresolved
primitive refines from 96 to at most 3072 interval bits before refusal. An ordered
104-bit batch path precedes that fallback. Execution binds evaluator sources,
Python and loaded math libraries, CPU dispatch fields, primitive backend, and
implementation manifests. Proposal libraries are single-threaded; OpenBLAS, OMP,
MKL and NumExpr thread counts are one. Changed bindings require requalification,
and replay must match the original feature-word commitments, not just versions.

The service, optimization, matched continuation and quality workers consumed
47,767, 35,258, 20,526 and 7,526 numerical CPU seconds, respectively. The matched
figure includes its failed-screen and diagnostic debit. Preparation costs and
all failures remain in the evidence record; worker and scheduler clocks overlap.
An unresolved historical 12,602-second reservation remains separate. These
figures are not a complete all-history compute total.

\paragraph{Inputs and responsibility.}
The author-side archive preserves outputs, losses, selections, failures and
1,322 source/protocol copies. Source release terms are unassigned; no external
review supplement is provided. Fresh reproduction from external inputs remains
unexecuted. Model cards declare Apache 2.0 for DistilGPT2 and MIT for GPT2; C4's
ODC-BY and Common Crawl conditions do not establish rights to every page.
Logical calibration deletion leaves external corpora and audit archives intact.
No content-level privacy or demographic audit was performed. Generative AI
assisted planning, code, analysis, literature checks and writing; the authors
remain responsible for all content.

\paragraph{Choosing an archive.}
The matched results favor optimized lossless storage when latency dominates.
Compression needs a storage budget that justifies longer repair time and a policy
for unavailable replay sources. These archives preserve a binary64 feature target;
a lower-precision target needs its own correctness and quality comparison.
Repair-state savings imply neither inference-memory nor serving-throughput gains.

\section{Conclusion}
Lossy calibration evidence can determine exact adaptive weight codes. Uniform
certificates, source replay, and canonical state construction make successful
repairs reusable for subsequent deletions. The measured operating point trades
modest storage savings for additional certification work and workload-dependent
source access.
\label{page:main-end}
\clearpage
\begin{thebibliography}{15}

\bibitem[Chee et~al.(2023)Chee, Cai, Kuleshov, and De~Sa]{quip}
Jerry Chee, Yaohui Cai, Volodymyr Kuleshov, and Christopher De~Sa. 2023.
\href{https://proceedings.neurips.cc/paper_files/paper/2023/hash/0df38cd13520747e1e64e5b123a78ef8-Abstract-Conference.html}{QuIP: 2-bit quantization of large language models with guarantees}.
In \emph{Advances in Neural Information Processing Systems}, volume 36, pages 4396--4429.

\bibitem[Chen et~al.(2026)Chen, Shabanzadeh, Crn\v{c}evi\'{c}, Hoefler, and Alistarh]{babai}
Jiale Chen, Yalda Shabanzadeh, Elvir Crn\v{c}evi\'{c}, Torsten Hoefler, and Dan Alistarh. 2026.
\href{https://arxiv.org/abs/2507.18553v4}{The geometry of LLM quantization: GPTQ as Babai's nearest plane algorithm}.
In \emph{International Conference on Learning Representations}. Author version: arXiv:2507.18553v4.

\bibitem[Frantar et~al.(2023)Frantar, Ashkboos, Hoefler, and Alistarh]{gptq}
Elias Frantar, Saleh Ashkboos, Torsten Hoefler, and Dan Alistarh.
2023. \emph{GPTQ: Accurate Post-Training Quantization for Generative Pre-trained
Transformers}. ICLR. Primary preprint:
\href{https://arxiv.org/abs/2210.17323v2}{arXiv:2210.17323, version 2}.

\bibitem[Ginart et~al.(2019)Ginart, Guan, Valiant, and Zou]{ginart}
Antonio A. Ginart, Melody Y. Guan, Gregory Valiant, and James Zou.
2019. \emph{Making AI Forget You: Data Deletion in Machine Learning}.
Advances in Neural Information Processing Systems 32.
\href{https://papers.nips.cc/paper_files/paper/2019/file/cb79f8fa58b91d3af6c9c991f63962d3-Paper.pdf}{Primary proceedings paper}.

\bibitem[Hugging Face(n.d.-a)]{distilgpt2}
Hugging Face. n.d.-a.
\href{https://huggingface.co/distilbert/distilgpt2/blob/2290a62682d06624634c1f46a6ad5be0f47f38aa/README.md}{DistilGPT2 model card}.
Model repository, pinned revision \nolinkurl{2290a62682d0}; accessed 28 October 2026.

\bibitem[Hugging Face(n.d.-b)]{gpt2card}
Hugging Face. n.d.-b.
\href{https://huggingface.co/openai-community/gpt2/blob/607a30d783dfa663caf39e06633721c8d4cfcd7e/README.md}{GPT-2 model card}.
Model repository, pinned revision \nolinkurl{607a30d783df}; accessed 28 October 2026.

\bibitem[Lang et~al.(2026)Lang, Helvits, and Shlezinger]{medu}
Natalie Lang, Alon Helvits, and Nir Shlezinger. 2026.
\href{https://doi.org/10.1109/ACCESS.2026.3663428}{Memory-efficient distributed unlearning}.
\emph{IEEE Access}, 14:24361--24378.

\bibitem[Merity et~al.(2017)Merity, Xiong, Bradbury, and Socher]{wikitext}
Stephen Merity, Caiming Xiong, James Bradbury, and Richard Socher. 2017.
\href{https://openreview.net/pdf?id=Byj72udxe}{Pointer sentinel mixture models}.
In \emph{International Conference on Learning Representations}.

\bibitem[Raffel et~al.(2020)Raffel, Shazeer, Roberts, Lee, Narang, Matena, Zhou, Li, and Liu]{c4}
Colin Raffel, Noam Shazeer, Adam Roberts, Katherine Lee, Sharan Narang,
Michael Matena, Yanqi Zhou, Wei Li, and Peter J. Liu.
2020. \emph{Exploring the Limits of Transfer Learning with a Unified Text-to-Text
Transformer}. Journal of Machine Learning Research, 21(140):1--67.
\href{https://jmlr.org/papers/v21/20-074.html}{Primary journal article}.

\bibitem[Rump(2010)]{rump}
Siegfried M. Rump. 2010.
\href{https://www.tuhh.de/ti3/rump/intlab/ActaNumerica2010.pdf}{Verification methods: Rigorous results using floating-point arithmetic}.
\emph{Acta Numerica}, 19:287--449.

\bibitem[Shewchuk(1997)]{shewchuk}
Jonathan Richard Shewchuk. 1997.
\href{https://people.eecs.berkeley.edu/~jrs/papers/robustr.pdf}{Adaptive precision floating-point arithmetic and fast robust geometric predicates}.
\emph{Discrete \& Computational Geometry}, 18(3):305--363.

\bibitem[Tavory et~al.(2026)Tavory, Gade, Sarig, Touitou, and Guy]{qss}
Ami Tavory, Shripad Gade, Tal Sarig, Noam Touitou, and Ido Guy.
2026. \emph{Exact Unlearning via Quantized Sufficient Statistics}.
\href{https://arxiv.org/html/2610.07197v1}{arXiv:2610.07197, version 1}.

\bibitem[Xiong et~al.(2023)Xiong, Li, Li, and Cai]{exactfun}
Zuobin Xiong, Wei Li, Yingshu Li, and Zhipeng Cai. 2023.
\href{https://doi.org/10.1109/ICDM58522.2023.00188}{Exact-Fun: An exact and efficient federated unlearning approach}.
In \emph{IEEE International Conference on Data Mining}, pages 1439--1444.

\bibitem[Zhang et~al.(2023)Zhang, Song, and Sun]{qebverif}
Yedi Zhang, Fu Song, and Jun Sun. 2023.
\href{https://doi.org/10.1007/978-3-031-37703-7_20}{QEBVerif: Quantization error bound verification of neural networks}.
In \emph{Computer Aided Verification}, volume 13965 of LNCS, pages 413--437. Springer.

\bibitem[Zhao et~al.(2026)Zhao, Wang, Chang, and He]{execcert}
Ziyu Zhao, Xinyu Wang, Xiao-Wen Chang, and Yixuan He.
2026. \emph{From Mathematical to Executable Certificates for Machine Unlearning}.
\href{https://arxiv.org/html/2610.02268v1}{arXiv:2610.02268, version 1}.
\end{thebibliography}
\end{document}
